% A LaTeX template for MSc Thesis submissions to 
% Politecnico di Milano (PoliMi) - School of Industrial and Information Engineering
%
% S. Bonetti, A. Gruttadauria, G. Mescolini, A. Zingaro
% e-mail: template-tesi-ingind@polimi.it
%
% Last Revision: October 2021
%
% Copyright 2021 Politecnico di Milano, Italy. NC-BY

\documentclass{Configuration_Files/PoliMi3i_thesis}

%------------------------------------------------------------------------------
%	REQUIRED PACKAGES AND  CONFIGURATIONS
%------------------------------------------------------------------------------

% CONFIGURATIONS
\usepackage{parskip} % For paragraph layout
\usepackage{setspace} % For using single or double spacing
\usepackage{emptypage} % To insert empty pages
\usepackage{multicol} % To write in multiple columns (executive summary)
\setlength\columnsep{15pt} % Column separation in executive summary
\setlength\parindent{0pt} % Indentation
\raggedbottom  

% PACKAGES FOR TITLES
\usepackage{titlesec}
% \titlespacing{\section}{left spacing}{before spacing}{after spacing}
\titlespacing{\section}{0pt}{3.3ex}{2ex}
\titlespacing{\subsection}{0pt}{3.3ex}{1.65ex}
\titlespacing{\subsubsection}{0pt}{3.3ex}{1ex}
\usepackage{color}

% PACKAGES FOR LANGUAGE AND FONT
\usepackage[english]{babel} % The document is in English  
\usepackage[utf8]{inputenc} % UTF8 encoding
\usepackage[T1]{fontenc} % Font encoding
\usepackage[11pt]{moresize} % Big fonts
\usepackage{enumitem}

% PACKAGES FOR IMAGES
\usepackage{graphicx}
\usepackage{svg}
\usepackage{transparent} % Enables transparent images
\usepackage{eso-pic} % For the background picture on the title page
\usepackage{subfig} % Numbered and caption subfigures using \subfloat.
\usepackage{tikz} % A package for high-quality hand-made figures.
\usetikzlibrary{}
\graphicspath{{./Images/}} % Directory of the images
\usepackage{caption} % Coloured captions
\usepackage{xcolor} % Coloured captions
\usepackage{amsthm,thmtools,xcolor} % Coloured "Theorem"
\usepackage{float}

% STANDARD MATH PACKAGES
\usepackage{amsmath}
\usepackage{amsthm}
\usepackage{amssymb}
\usepackage{amsfonts}
\usepackage{bm}
\usepackage[overload]{empheq} % For braced-style systems of equations.
\usepackage{fix-cm} % To override original LaTeX restrictions on sizes

% PACKAGES FOR TABLES
\usepackage{tabularx}
\usepackage{longtable} % Tables that can span several pages
\usepackage{colortbl}
\usepackage{geometry}
\usepackage{array}      % per newcolumntype e \arraybackslash
\usepackage{booktabs}   % per \toprule \midrule \bottomrule
\usepackage{pifont}     % per i segni di spunta/croce
\newcommand{\cmark}{\ding{51}}
\newcommand{\xmark}{\hspace*{0.25cm}}
\newcolumntype{R}[1]{>{\raggedleft\arraybackslash}p{#1}} 


% PACKAGES FOR ALGORITHMS (PSEUDO-CODE)
\usepackage{algorithm}
\usepackage{algorithmic}

% PACKAGES FOR REFERENCES & BIBLIOGRAPHY
\usepackage[colorlinks=true,linkcolor=black,anchorcolor=black,citecolor=black,filecolor=black,menucolor=black,runcolor=black,urlcolor=black]{hyperref} % Adds clickable links at references
\usepackage{cleveref}
\usepackage[square, numbers, sort&compress]{natbib} % Square brackets, citing references with numbers, citations sorted by appearance in the text and compressed
\bibliographystyle{abbrvnat} % You may use a different style adapted to your field

% OTHER PACKAGES
\usepackage{pdfpages} % To include a pdf file
\usepackage{afterpage}
\usepackage{lipsum} % DUMMY PACKAGE
\usepackage{fancyhdr} % For the headers
\fancyhf{}

% Input of configuration file. Do not change config.tex file unless you really know what you are doing. 
% Define blue color typical of polimi
\definecolor{bluepoli}{cmyk}{0.4,0.1,0,0.4}

% Custom theorem environments
\declaretheoremstyle[
  headfont=\color{bluepoli}\normalfont\bfseries,
  bodyfont=\color{black}\normalfont\itshape,
]{colored}

% Set-up caption colors
\captionsetup[figure]{labelfont={color=bluepoli}} % Set colour of the captions
\captionsetup[table]{labelfont={color=bluepoli}} % Set colour of the captions
\captionsetup[algorithm]{labelfont={color=bluepoli}} % Set colour of the captions

\theoremstyle{colored}
\newtheorem{theorem}{Theorem}[chapter]
\newtheorem{proposition}{Proposition}[chapter]

% Enhances the features of the standard "table" and "tabular" environments.
\newcommand\T{\rule{0pt}{2.6ex}}
\newcommand\B{\rule[-1.2ex]{0pt}{0pt}}

% Pseudo-code algorithm descriptions.
\newcounter{algsubstate}
\renewcommand{\thealgsubstate}{\alph{algsubstate}}
\newenvironment{algsubstates}
  {\setcounter{algsubstate}{0}%
   \renewcommand{\STATE}{%
     \stepcounter{algsubstate}%
     \Statex {\small\thealgsubstate:}\space}}
  {}

% New font size
\newcommand\numfontsize{\@setfontsize\Huge{200}{60}}

% Title format: chapter
\titleformat{\chapter}[hang]{
\fontsize{50}{20}\selectfont\bfseries\filright}{\textcolor{bluepoli} \thechapter\hsp\hspace{2mm}\textcolor{bluepoli}{|   }\hsp}{0pt}{\huge\bfseries \textcolor{bluepoli}
}

% Title format: section
\titleformat{\section}
{\color{bluepoli}\normalfont\Large\bfseries}
{\color{bluepoli}\thesection.}{1em}{}

% Title format: subsection
\titleformat{\subsection}
{\color{bluepoli}\normalfont\large\bfseries}
{\color{bluepoli}\thesubsection.}{1em}{}

% Title format: subsubsection
\titleformat{\subsubsection}
{\color{bluepoli}\normalfont\large\bfseries}
{\color{bluepoli}\thesubsubsection.}{1em}{}

% Shortening for setting no horizontal-spacing
\newcommand{\hsp}{\hspace{0pt}}

\makeatletter
% Renewcommand: cleardoublepage including the background pic
\renewcommand*\cleardoublepage{%
  \clearpage\if@twoside\ifodd\c@page\else
  \null
  \AddToShipoutPicture*{\BackgroundPic}
  \thispagestyle{empty}%
  \newpage
  \if@twocolumn\hbox{}\newpage\fi\fi\fi}
\makeatother

%For correctly numbering algorithms
\numberwithin{algorithm}{chapter}

%----------------------------------------------------------------------------
%	NEW COMMANDS DEFINED
%----------------------------------------------------------------------------

% EXAMPLES OF NEW COMMANDS
\newcommand{\bea}{\begin{eqnarray}} % Shortcut for equation arrays
\newcommand{\eea}{\end{eqnarray}}
\newcommand{\e}[1]{\times 10^{#1}}  % Powers of 10 notation

%----------------------------------------------------------------------------
%	ADD YOUR PACKAGES (be careful of package interaction)
%----------------------------------------------------------------------------

%----------------------------------------------------------------------------
%	ADD YOUR DEFINITIONS AND COMMANDS (be careful of existing commands)
%----------------------------------------------------------------------------

%----------------------------------------------------------------------------
%	BEGIN OF YOUR DOCUMENT
%----------------------------------------------------------------------------

\begin{document}

\fancypagestyle{plain}{%
\fancyhf{} % Clear all header and footer fields
\fancyhead[RO,RE]{\thepage} %RO=right odd, RE=right even
\renewcommand{\headrulewidth}{0pt}
\renewcommand{\footrulewidth}{0pt}}

%----------------------------------------------------------------------------
%	TITLE PAGE
%----------------------------------------------------------------------------

\pagestyle{empty} % No page numbers
\frontmatter % Use roman page numbering style (i, ii, iii, iv...) for the preamble pages

\puttitle{
	title=Software Engineering 2 - DD, % Title of the thesis
	name=Simone Controguerra,
        name=Matteo Delogu,
            name=Gianluca Di Paola,% Author Name and Surname % Student ID number (numero di matricola)
	advisor= Prof. Elisabetta Di Nitto, % Supervisor name
	 % Co-Supervisor name, remove this line if there is none
	academicyear={2025-2026}  % Academic Year
} 

%----------------------------------------------------------------------------
%	PREAMBLE PAGES: ABSTRACT (inglese e italiano), EXECUTIVE SUMMARY
%----------------------------------------------------------------------------
\startpreamble
\setcounter{page}{1} % Set page counter to 1

%----------------------------------------------------------------------------
%	LIST OF CONTENTS/FIGURES/TABLES/SYMBOLS
%----------------------------------------------------------------------------

% TABLE OF CONTENTS
\thispagestyle{empty}
\tableofcontents % Table of contents 
\thispagestyle{empty}
\cleardoublepage

%-------------------------------------------------------------------------
%	THESIS MAIN TEXT
%-------------------------------------------------------------------------
% In the main text of your thesis you can write the chapters in two different ways:
%
%(1) As presented in this template you can write:
%    \chapter{Title of the chapter}
%    *body of the chapter*
%
%(2) You can write your chapter in a separated .tex file and then include it in the main file with the following command:
%    \chapter{Title of the chapter}
%    \input{chapter_file.tex}
%
% Especially for long thesis, we recommend you the second option.

\addtocontents{toc}{\vspace{2em}} % Add a gap in the Contents, for aesthetics
\mainmatter % Begin numeric (1,2,3...) page numbering

% --------------------------------------------------------------------------
% NUMBERED CHAPTERS % Regular chapters following
% --------------------------------------------------------------------------

\chapter{Introduction}

\section{Purpose}
Nowadays self-care activities have become really popular such as keeping themselves fit, in particular a lot of people ride their bike either to stay healthy or to adopt an eco-friendly means of transport. Although, streets cannot be the safest place for bikes since cars can run very fast there. The solution is to ride on designated bike paths or safe roads where cars are few and they cannot run too fast.

Best Bike Paths is a service that offers an innovative solution to bikers that want to ride safely. Its purpose is whether to give the opportunity to its users to record their personal trips and store them as an archive of sport activities, enriched with different kinds of related information, or to provide suggestions for safe and efficient bike paths to reach a given destination, highlighting the better ones. Best Bike Paths is a service with community features, that aims to build a network of bike paths for everyone.
\section{Scope}
Best Bike Paths is a system that collects and processes cycling data through the coordination of software components. Its scope includes managing registered users and their personal trips, acquiring data from mobile devices' sensors, and producing processed information such as metrics and automatically inferred bike paths. The system maintains a shared collection of bike paths, updated through user reports, ensuring that public streets remain accurate and reliable. Best Bike Path provides a search service that allows the users to find suitable paths for their destination.
\section{Definitions, Acronyms and Abbreviations}
\begin{itemize}

    \item \textbf{BBP (Best Bike Paths):} The name of the system.
    \item \textbf{Bike Path:} A route specifically designated or suitable for bicycle travel, which can be either manually defined by users or automatically generated by the system. It is formally defined as a set consisting of: a Starting Point, a Finishing Point, and an ordered sequence of traversable Streets (see definition of Street).
    \item \textbf{Trip:} A standalone, recorded instance of a user traveling from one point to another, which includes sensor data, metrics, and potentially weather information.
    \item \textbf{Street:} A segment of road suitable for cycling (factors taken into account include average bike speed limit, roadway width, traffic volume, dedicated bike lanes, etc.) that forms part of a bike path, with attributes such as status and presence of irregularities (e.g., potholes).
    \item \textbf{Navigation System:} An external service or API used by the system to evaluate road safety for bikers, considering factors such as bike track availability and traffic conditions.
    \item \textbf{Weather-forecast System:} An external service or API used by the system to provide weather information associated with a trip.
    \item \textbf{GPS (Global Positioning System):} A satellite-based navigation system that provides location and time information, in all weather conditions, anywhere on Earth.
    \item \textbf{IMU (Inertial Measurement Unit) Sensors:} Electronic devices that measure and report a physical object's specific force, angular rate and orientation, using a combination of both accelerometer and gyroscope.
    \item \textbf{API (Application Programming Interface):} A set of protocols and tools that enable different software applications to communicate with each other.
    \item \textbf{JSON (JavaScript Object Notation):} A lightweight data-interchange format that is easily readable and writable by both humans and machines.
    \item \textbf{UI (User Interface):} The interface through which users interact with the platform, encompassing buttons, screens, and navigation elements.
    \item \textbf{TTL (time-to-live ):} Maximum duration for a cache entry to be considered valid.
    \item \textbf{Redundancy:} The inclusion of supplementary components or systems to ensure reliability and availability in the event of a primary component failure.
    \item \textbf{Scalability:} The system's capacity to manage a growing number of users or increased load by incorporating additional resources without significant degradation in performance.
    \item \textbf{Uptime:} The period during which a system is fully operational and accessible to users without interruption.
    \item \textbf{Backup:} The process of duplicating and archiving data to ensure restorability in the event of loss or damage to the original data.
    
    %dovremmo aggiungere la definizione di cache?

\end{itemize}
\section{Revision History}
\begin{itemize}
    \item \textbf{v1.0} First version of the document, revised on xxth of December 2025.
\end{itemize}
\section{Reference Documents}
Professor Elisabetta Di Nitto, \textit{Slides of the course "Software Engineering 2"} at PoliMi.
\\ 
Polytechnic University of Milan, \textit{Specifications of the DD assignment for the Software Engineering II course}.
\section{Document Structure}
\begin{enumerate}
    \item \textbf{Introduction:} it provides an overview about the project.
    \item \textbf{Architectural Design:} 
    \item \textbf{User Interface Design:} 
    \item \textbf{Requirements traceability:} 
    \item \textbf{Implementation, Integration and Test plan:} 
    \item \textbf{Effort spent:} it indicates how the workload has been distributed among the contributors.
    \item \textbf{References:} it includes references to documents and tools that have been used in the writing of the document. 
\end{enumerate}
%------------------------------------------------------
\chapter{Architectural Design}

\section{Overview: High-level components and their interaction}

\section{Component View}

\section{Deployment View}

\section{Runtime View}

\section{Component Interfaces}

\section{Selected Architectural Styles and Patterns}
\textit{TODO: Gianluca}

\subsection{Load Balancing}
%%another section may be needed

%------------------------------------------------------
\chapter{User Interface Design}
The details of the UI have already been analyzed and showed in the User Interface section of the RASD document.

%------------------------------------------------------
\chapter{Requirements traceability}
\textit{TODO: Gianluca}

%------------------------------------------------------
\chapter{Implementation, Integration and Test plan}

\section{Component Priority}

\section{Integration and testing Plan}

\section{System Testing}

%------------------------------------------------------
\chapter{Effort spent}

\chapter{References}


%------------------------------------------------------
\listoffigures

\listoftables

%------------------------------------------------------
%------------------------------------------------------
\cleardoublepage
\end{document}
